\documentclass[10pt,a4paper]{article}
\usepackage[T1]{fontenc}
\usepackage{lmodern}
\usepackage[margin=2.2cm]{geometry}
\usepackage{microtype}
\usepackage{parskip}
\usepackage[hidelinks]{hyperref}
\pagestyle{empty}

\begin{document}

\noindent
{\large Luca Cirfeta}\\
Independent Researcher, Rome, Italy\\
\href{mailto:lucacirfeta@gmail.com}{lucacirfeta@gmail.com}\\
ORCID: \href{https://orcid.org/0009-0000-1235-3186}{0009-0000-1235-3186}

\vspace{0.7em}
\noindent
\today

\vspace{0.7em}
\noindent
The Editor\\
\emph{Classical and Quantum Gravity}\\
IOP Publishing

\vspace{0.9em}
\noindent Dear Editor,

I submit the Research Paper ``Detector-Dependent Domain Shift and Failure-Mode
Validation in Unsupervised LIGO Transient-Noise Characterization with DANTE''
for consideration in \emph{Classical and Quantum Gravity}.

An earlier manuscript (CQG-116729) was declined before external review.  This
submission is a scientific and editorial reconstruction, not a cosmetic
resubmission: its title, claim, population, results, structure, and reference
survey have changed.

The Domain-Adaptive Network for Transient Evaluation (DANTE) addresses how to
validate unsupervised anomaly searches when detector noise changes between runs
and target-run adaptation can absorb recurrent anomalies. Its central claim is
methodological rather than a candidate announcement: representation, reference
and calibration populations, and physical nulls must be explicit and separately
checkable.

All class-dependent results use a coherent Q64/Q64 native index and independent
background calibration.  The detector-aware catalogue contains 10,429 O4a
candidates; correcting the former incoherent representation changes 4,676 of
10,372 paired dispositions.  Direct controls show detector-dependent domain
shift and morphology-dependent known-glitch separation.  Replicated
sensitivity and negative-control studies measure stability, native-index
absorption, and a conditional low-Q blind region.

Physical controls avoid discovery language. A non-exchangeable H1--L1 screen
finds 13/8,806 values above its threshold but assigns neither a tail $p$-value
nor physical meaning. The family-wise auxiliary-channel endpoint does not
resolve class enrichment (2/26 versus 7/93; Fisher $p=1.000$), and two catalogue
overlaps are consistent with circular shifts ($p=0.651$). The former discovery
and rate-limit interpretations are withdrawn. The retained L1 candidate is
reported as an unclassified local transient, not as a new glitch class.

The manuscript includes explicit limitations, disclosures, a validation
matrix, and requirements for another observing run.  Public data are cited
through GWOSC; DANTE 3.7.0 is archived at DOI
10.5281/zenodo.21912589, and the separate v6 evidence dataset at DOI
10.5281/zenodo.21925453.

The relationship to public preprints is explicit. arXiv:2607.18136 reports the
earlier O4a application and the interpretations reassessed here.
arXiv:2608.15166 is the shorter representation-coherent reanalysis underlying
this submission. The submitted article is a self-contained and substantially
expanded journal treatment: it provides the complete method, broader related-
work comparison, statistical contract, additional controls, numerical
validation matrix, reproducibility details, and requirements for application
to another observing run.

The work has not been published in a peer-reviewed journal and is not under
consideration elsewhere. I am the sole author and accept responsibility for
the analysis and manuscript. Classical and Quantum Gravity's broad readership
of gravitational-wave experimentalists and theorists makes it the appropriate
venue. Referees with expertise in detector characterization and empirical
search calibration would be particularly appropriate.

Thank you for considering this substantially revised work.

\vspace{0.35em}
\noindent Yours sincerely,
\nopagebreak\\[0.35em]
Luca Cirfeta

\end{document}
